\documentclass[a4paper,12pt]{article}
\usepackage{iftex}
\ifPDFTeX
  \usepackage[T2A]{fontenc}
  \usepackage[utf8]{inputenc}
  \usepackage[russian]{babel}
\else
  \usepackage{fontspec}
  \setmainfont{DejaVu Sans}
  \usepackage{polyglossia}
  \setdefaultlanguage{russian}
\fi
\usepackage{amsmath,amssymb,amsthm,mathtools}
\usepackage{icomma}
\usepackage[a4paper,margin=1.65cm,headheight=16pt]{geometry}
\usepackage{microtype}
\usepackage{graphicx}
\usepackage{tikz}
\usetikzlibrary{positioning,arrows.meta,calc}
\usepackage{tkz-euclide}
\usepackage{pgfplots}
\pgfplotsset{compat=1.18}
\usepackage{tabularx,booktabs,longtable}
\usepackage{enumitem}
\usepackage{hyperref}
\usepackage{parskip}
\usepackage{fancyhdr}
\usepackage{setspace}
\usepackage{xcolor}
\usepackage[most]{tcolorbox}
\usepackage{multicol}

\definecolor{accent}{HTML}{287F7A}
\definecolor{darkaccent}{HTML}{1C3151}
\definecolor{textgray}{HTML}{333333}
\definecolor{softteal}{HTML}{E4F3EF}
\definecolor{softcoral}{HTML}{FBE4DC}
\definecolor{coral}{HTML}{E77B5E}
\definecolor{softgold}{HTML}{F8EDCF}
\definecolor{softblue}{HTML}{E8EEF6}

\hypersetup{colorlinks=true,linkcolor=darkaccent,urlcolor=accent}
\setlist[itemize]{leftmargin=1.5em,itemsep=0.25em,topsep=0.35em}
\setlist[enumerate]{leftmargin=1.8em,itemsep=0.35em,topsep=0.35em}
\pagestyle{fancy}
\fancyhf{}
\lhead{Комплексное повторение ЕГЭ}
\rhead{Николь Саркисьянц}
\cfoot{\thepage}
\renewcommand{\headrulewidth}{0.4pt}
\renewcommand{\today}{21 июля 2026 г.}

\newtcolorbox{keybox}[1][]{enhanced,breakable,colback=softteal,colframe=accent,boxrule=0.8pt,arc=3mm,left=4mm,right=4mm,top=3mm,bottom=3mm,title=#1,fonttitle=\bfseries,coltitle=darkaccent}
\newtcolorbox{warnbox}[1][]{enhanced,breakable,colback=softcoral,colframe=coral,boxrule=0.8pt,arc=3mm,left=4mm,right=4mm,top=3mm,bottom=3mm,title=#1,fonttitle=\bfseries,coltitle=darkaccent}
\newtcolorbox{examplebox}[1][]{enhanced,breakable,colback=softblue,colframe=darkaccent,boxrule=0.8pt,arc=3mm,left=4mm,right=4mm,top=3mm,bottom=3mm,title=#1,fonttitle=\bfseries,coltitle=darkaccent}
\newcommand{\R}{\mathbb{R}}

\begin{document}
\begin{titlepage}
\centering
\vspace*{2.1cm}
{\color{accent}\Large\bfseries ЧЕК-ЛИСТ\par}
\vspace{0.45cm}
{\color{darkaccent}\fontsize{29}{34}\selectfont\bfseries Комплексное повторение\\тестовой части ЕГЭ\par}
\vspace{0.8cm}
{\large Финансы, геометрия, векторы, вероятность, степени, логарифмы и текстовые задачи\par}
\vfill
{\Large\bfseries Лёвин Артём Александрович\\[0.2cm]эксклюзивно для Николь Саркисьянц\par}
\vspace{0.8cm}
{\large \today\par}
\vfill
\begin{tikzpicture}[remember picture,overlay]
  \draw[accent,line width=2.2pt] (-7,-1.1) .. controls (-4.6,1.1) and (-2,-2.0) .. (0.1,0.05)
    .. controls (2.2,2.0) and (4.8,-0.9) .. (7,0.75);
  \fill[accent!35] (-6.0,-0.30) circle (3pt);
  \fill[accent!60] (0.15,-0.12) circle (3pt);
  \fill[accent!85] (6.15,0.20) circle (3pt);
\end{tikzpicture}
\end{titlepage}

\section*{1. Главная стратегия повторительного варианта}
\begin{keybox}[Единый алгоритм]
\begin{enumerate}
\item Определите тип задачи и выпишите известные данные.
\item Выберите формулу или модель, которая непосредственно связывает данные с вопросом.
\item Переведите слова условия в математические знаки: ``сняли'' означает минус, ``больше'' --- строгую границу, ``не менее'' --- включённую границу.
\item Выполните преобразования без скачков: каждый процент считайте от суммы на начало соответствующего периода.
\item Проверьте знак, единицы измерения, область допустимых значений и смысл ответа.
\end{enumerate}
\end{keybox}

\begin{center}
\begin{tikzpicture}[node distance=7mm and 8mm,>=Stealth,every node/.style={font=\small,align=center,rounded corners=2.5mm,minimum height=9mm}]
\node[draw,fill=softgold,text width=3.2cm] (a) {Прочитать условие};
\node[draw,fill=softteal,text width=3.2cm,right=of a] (b) {Выбрать модель};
\node[draw,fill=softblue,text width=3.2cm,right=of b] (c) {Вычислить};
\node[draw,fill=softcoral,text width=3.2cm,right=of c] (d) {Проверить смысл};
\draw[->,thick,accent] (a)--(b);
\draw[->,thick,accent] (b)--(c);
\draw[->,thick,accent] (c)--(d);
\end{tikzpicture}
\end{center}

\section*{2. Вклады и проценты}
\begin{keybox}[Два принципа]
\begin{itemize}
\item Если деньги снимают со счёта, операция записывается со знаком минус.
\item Процент очередного года начисляется на сумму, которая находится на счёте в начале этого года.
\end{itemize}
\end{keybox}

\begin{examplebox}[Ставки по годам]
После двух лет на счёте:
\[
S_2=1{,}1^2S.
\]
После третьего года:
\[
S_3=1{,}1^2S+1{,}1^2S\cdot\frac{x}{100}
=1{,}1^2S\left(1+\frac{x}{100}\right).
\]
Степень $3$ у множителя $1{,}1$ здесь не появляется: ставка третьего года равна $x\%$, а не $10\%$.
\end{examplebox}

\begin{warnbox}[Типичная ошибка]
Нельзя считать третий процент от первоначальной суммы $S$. База начисления --- текущая сумма $1{,}1^2S$.
\end{warnbox}

\section*{3. Планиметрия: равнобедренный треугольник}
\begin{center}
\begin{tikzpicture}[scale=0.9]
\tkzDefPoint(-3,0){A}
\tkzDefPoint(3,0){B}
\tkzDefPoint(0,4){C}
\tkzDefPoint(4.4,0){D}
\tkzDrawPolygon[thick,color=darkaccent](A,B,C)
\tkzDrawSegment[thick,color=coral](B,D)
\tkzMarkSegments[mark=|,color=accent](A,C B,C)
\tkzMarkAngle[size=0.65,color=coral,mark=none](C,B,D)
\tkzLabelAngle[pos=1.0,color=coral](C,B,D){$107^\circ$}
\tkzLabelPoints[below](A,B,D)
\tkzLabelPoints[above](C)
\end{tikzpicture}
\end{center}

Если $AC=BC$, то углы при основании равны. Внешний угол при $B$ равен $107^\circ$, поэтому внутренний угол
\[
\angle CBA=180^\circ-107^\circ=73^\circ.
\]
Следовательно, $\angle CAB=73^\circ$, а
\[
\angle ACB=180^\circ-73^\circ-73^\circ=34^\circ.
\]

\section*{4. Векторы}
\begin{keybox}[Координаты и скалярное произведение]
Если $\vec a=(a_x;a_y)$ и $\vec b=(b_x;b_y)$, то
\[
\vec a\cdot\vec b=a_xb_x+a_yb_y.
\]
Одновременно
\[
\vec a\cdot\vec b=|\vec a|\,|\vec b|\cos\varphi,
\qquad
|\vec a|=\sqrt{a_x^2+a_y^2}.
\]
Отсюда
\[
\cos\varphi=
\frac{a_xb_x+a_yb_y}
{\sqrt{a_x^2+a_y^2}\sqrt{b_x^2+b_y^2}}.
\]
\end{keybox}

\begin{examplebox}[По клеткам координатной плоскости]
Пусть $\vec a=(2;3)$, $\vec b=(2;-1)$. Тогда
\[
\vec a\cdot\vec b=2\cdot2+3\cdot(-1)=1.
\]
Знак координаты по вертикали определяется направлением: движение вниз даёт отрицательную координату.
\end{examplebox}

\begin{warnbox}[Выбор формулы]
Формула $|\vec a||\vec b|\cos\varphi$ неудобна, если неизвестны длины и угол. При известных координатах начинайте с координатной формулы.
\end{warnbox}

\section*{5. Вероятность: вложенные события}
Пусть $X$ --- число решённых задач. Если
\[
P(X>4)=0{,}73,\qquad P(X>3)=0{,}86,
\]
то событие $\{X>4\}$ вложено в $\{X>3\}$, поэтому
\[
P(X=4)=P(X>3)-P(X>4)=0{,}86-0{,}73=0{,}13.
\]

\begin{center}
\begin{tikzpicture}[x=1.15cm,y=1cm]
\draw[->,thick] (0,0)--(10.2,0);
\foreach \x in {0,...,9} {\draw (\x,0.12)--(\x,-0.12) node[below=3pt]{\small \x};}
\draw[line width=7pt,coral!55,line cap=round] (4,0.55)--(9.6,0.55);
\draw[line width=7pt,accent!70,line cap=round] (5,0.9)--(9.6,0.9);
\node[anchor=west] at (0,0.9) {\small $X>4$};
\node[anchor=west] at (0,0.55) {\small $X>3$};
\end{tikzpicture}
\end{center}

\begin{keybox}[Перевод слов]
\begin{tabularx}{\linewidth}{@{}lX@{}}
\toprule
Формулировка & Для целого $X$ \\
\midrule
больше $8$ & $X\ge 9$ \\
не менее $8$ & $X\ge 8$ \\
меньше $8$ & $X\le 7$ \\
не более $8$ & $X\le 8$ \\
\bottomrule
\end{tabularx}
\end{keybox}

\section*{6. Степени и логарифмы}
\begin{examplebox}[Показательное уравнение]
Если
\[
3^{x-8}=\frac1{81}=3^{-4},
\]
то основания одинаковы и положительны, поэтому
\[
x-8=-4,\qquad x=4.
\]
\end{examplebox}

\begin{examplebox}[Разность логарифмов]
\[
\log_2 56-\log_2 7
=\log_2\frac{56}{7}
=\log_2 8=3.
\]
Объединять логарифмы можно при одинаковых основаниях и отсутствии дополнительных коэффициентов перед ними.
\end{examplebox}

\section*{7. Прикладные формулы и текстовые задачи}
\begin{warnbox}[Буква не равна цифре]
В формуле внимательно различайте обозначения. Например, $L$ --- пройденный путь, а не цифра $1$. Перед подстановкой подпишите смысл каждой переменной и её единицы.
\end{warnbox}

\subsection*{Задача на движение}
Два велосипедиста проходят по $190$ км. Быстрый едет со скоростью $x$ км/ч, медленный --- $(x-9)$ км/ч. Первый прибывает на $9$ часов раньше:
\[
\frac{190}{x-9}-\frac{190}{x}=9,\qquad x>9.
\]
После умножения на $x(x-9)$:
\[
190x-190(x-9)=9x(x-9),
\]
\[
x^2-9x-190=0.
\]
\[
D=81+760=841=29^2,\qquad x=\frac{9+29}{2}=19.
\]
Отрицательный корень не удовлетворяет смыслу задачи. Ответ: $19$ км/ч.

\begin{center}
\renewcommand{\arraystretch}{1.35}
\begin{tabularx}{0.84\linewidth}{@{}lXXX@{}}
\toprule
& Скорость & Время & Путь \\
\midrule
Первый & $x$ & $\dfrac{190}{x}$ & $190$ \\
Второй & $x-9$ & $\dfrac{190}{x-9}$ & $190$ \\
\bottomrule
\end{tabularx}
\end{center}

\section*{8. Как извлекать корень из большого квадрата}
\begin{keybox}[Оценка по десяткам и последней цифре]
\begin{enumerate}
\item Найдите соседние круглые квадраты: например, $70^2$ и $80^2$.
\item Определите десяток искомого корня.
\item По последней цифре подберите две возможные цифры единиц.
\item Сравните число с ближайшим круглым квадратом и проверьте кандидата.
\end{enumerate}
\end{keybox}

\begin{examplebox}[Пример]
\[
70^2=4900<5329<6400=80^2.
\]
Корень имеет вид $7\Box$. Число оканчивается на $9$, поэтому кандидаты --- $73$ и $77$. Число ближе к $4900$, проверяем:
\[
73^2=5329,\qquad \sqrt{5329}=73.
\]
Аналогично:
\[
\sqrt{7569}=87,\qquad \sqrt{1764}=42.
\]
\end{examplebox}

\section*{9. Фрагмент рабочей таблицы занятия}
\IfFileExists{../table.png}{
\begin{center}
\includegraphics[width=\linewidth,height=0.52\textheight,keepaspectratio]{../table.png}
\end{center}
}{
\begin{center}
\begin{tabularx}{\linewidth}{@{}lXX@{}}
\toprule
Блок & Главный инструмент & Проверка \\
\midrule
Финансы & коэффициент роста и знаки операций & процент от текущей суммы \\
Векторы & координатная формула & знак вертикальной координаты \\
Вероятность & вложенные события & разность большего и меньшего события \\
Движение & таблица $v,t,s$ & положительность скорости \\
Корни & оценка между квадратами & возведение кандидата в квадрат \\
\bottomrule
\end{tabularx}
\end{center}
}

\section*{10. Типичные ошибки}
\begin{warnbox}
\begin{itemize}
\item Снятие денег записано со знаком плюс.
\item Процент вычислен от первоначальной, а не от текущей суммы.
\item Внешний угол принят за внутренний.
\item Вектор, направленный вниз, получил положительную $y$-координату.
\item Вероятности вычитаются без проверки вложенности событий.
\item Слова ``больше'' и ``не менее'' считаются равнозначными.
\item Буква $L$ прочитана как цифра $1$.
\item В задаче на движение не учтено условие $x>9$.
\item Квадратный корень угадывается без итоговой проверки.
\end{itemize}
\end{warnbox}

\section*{11. Итоговый чек-лист}
\begin{itemize}[label=$\square$]
\item Я различаю пополнение и снятие по знаку.
\item Я начисляю процент на сумму начала текущего периода.
\item Я умею перейти от внешнего угла к внутреннему.
\item Я читаю координаты вектора по горизонтальному и вертикальному смещению.
\item Я знаю обе формулы скалярного произведения и длину вектора.
\item Я вычитаю вероятности только для вложенных событий.
\item Я точно перевожу формулировки ``больше'', ``не менее'', ``меньше'', ``не более''.
\item Я свожу степени к одинаковому основанию.
\item Я применяю свойства логарифмов только при выполнении условий.
\item Я подписываю смысл букв в прикладной формуле.
\item Я строю таблицу скорость--время--путь и проверяю ограничения.
\item Я извлекаю корень оценкой и проверяю кандидата.
\end{itemize}

\clearpage
\section*{Тренировочный блок}
\subsection*{Средний уровень}
\begin{enumerate}
\item В равнобедренном треугольнике боковые стороны равны, внешний угол при основании равен $116^\circ$. Найдите угол при вершине.
\item Даны $\vec a=(3;-2)$ и $\vec b=(1;4)$. Найдите $\vec a\cdot\vec b$.
\item $P(X>5)=0{,}62$, $P(X>4)=0{,}79$. Найдите $P(X=5)$.
\end{enumerate}

\subsection*{Уровень выше среднего}
\begin{enumerate}[resume]
\item Найдите косинус угла между $\vec a=(1;2)$ и $\vec b=(2;-1)$.
\item Решите уравнение $5^{x-2}=\dfrac1{125}$.
\item Вычислите $\log_3 54-\log_3 2$.
\item Два велосипедиста проходят $240$ км. Скорость первого на $10$ км/ч больше, и он прибывает на $2$ часа раньше. Найдите скорость первого.
\item На счёте было $100\,000$ рублей. После начисления $8\%$ сняли $5\,000$ рублей. Сколько осталось?
\item Найдите $\sqrt{6241}$ методом оценки.
\end{enumerate}

\subsection*{Сложный уровень}
\begin{enumerate}[resume]
\item Вклад $S$ два года рос на $10\%$ ежегодно, а в третий год --- на $x\%$. После третьего года сумма стала равна $1{,}452S$. Найдите $x$.
\end{enumerate}

\clearpage
\section*{Ответы}
\renewcommand{\arraystretch}{1.45}
\begin{longtable}{@{}cl@{}}
\toprule
№ & Ответ \\
\midrule
1 & $52^\circ$ \\
2 & $-5$ \\
3 & $0{,}17$ \\
4 & $0$ \\
5 & $-1$ \\
6 & $3$ \\
7 & $40$ км/ч \\
8 & $103\,000$ рублей \\
9 & $79$ \\
10 & $20\%$ \\
\bottomrule
\end{longtable}

\end{document}
